%%%% CAPÍTULO 5 - CONCLUSÕES E PERSPECTIVAS
%%
\chapter{Considerações finais}\label{cap:conclusoes}

Neste trabalho, foi evidenciada a relevância dos estudos na área de arquitetura de computadores, consolidada como um dos pilares dos grandes avanços tecnológicos modernos, bem como o impacto positivo que as arquiteturas de código aberto têm exercido nesse ecossistema. Diante disso, estabeleceu-se o objetivo de projetar e implementar um núcleo processador baseado na ISA RISC-V em uma plataforma FPGA. Para a fundamentação teórica deste projeto, realizou-se uma revisão detalhada de conceitos clássicos, como a Arquitetura de Harvard e seus blocos estruturais essenciais: a Unidade Lógica e Aritmética (ALU), a Unidade de Controle (CU), os sistemas de Entrada e Saída (I/O) e as memórias distintas de dados e de instruções.

A pesquisa delimitou o conjunto de instruções a ser adotado, analisando as extensões da ISA e definindo, por critérios de viabilidade e complexidade, a implementação do conjunto básico inteiro de 32 bits (RV32I). Na proposta do trabalho, foram definidas as especificações de recursos de hardware e software, destacando-se a placa de desenvolvimento Nexys A7-100T e a suíte AMD Vivado. Complementando o planejamento, estruturou-se uma metodologia sistemática dividida em sete etapas estratégicas, além de um mapeamento inicial para a organização modular e hierárquica do fluxo de dados e das rotinas de teste.

Por fim, reconhece-se que a execução física deste projeto apresenta desafios técnicos complexos, especialmente no que tange à restrição de tempo para a sua finalização. Contudo, o planejamento e a divisão modular estabelecidos neste documento fornecem as bases para viabilizar a conclusão do sistema dentro dos prazos estipulados. Desse modo, os próximos passos deste trabalho consistirão no início da pesquisa detalhada sobre a modelagem dos componentes em HDL para FPGA, seguida da implementação e simulação dos módulos lógicos do processador de forma isolada e posteriormente uni-los.